% SCP10 arXiv1001.3807v3, Lemma 6.14: finite-coordinate bridge only.
\begin{definition}[Open coefficients on labelled bonds]
    \label{def:peps_labelled_open_coefficient}
    \lean{TNLean.PEPS.DependentBondNetwork.RegionBoundaryEndpoint,
        TNLean.PEPS.DependentBondNetwork.RegionSiteConfig,
        TNLean.PEPS.DependentBondNetwork.joinRegionConfig,
        TNLean.PEPS.DependentBondNetwork.openBondFactor,
        TNLean.PEPS.DependentBondNetwork.openCoefficient,
        TNLean.PEPS.DependentBondNetwork.openBoundaryContraction,
        TNLean.PEPS.DependentBondNetwork.deltaExteriorConfig,
        TNLean.PEPS.DependentBondNetwork.deltaCompletedTensor,
        TNLean.PEPS.DependentBondNetwork.internalBondMatrices}
    \notready
    For a region $R$ in a finite directed multigraph, assign an independent
    virtual index to every incidence at every vertex of $R$. Fix the indices
    on region-side incidences of crossing bonds to boundary data $\theta$.
    Multiply the region site coefficients and the prescribed matrix entries
    on internal bonds, and sum the remaining indices. Pair this coefficient
    with any joint boundary tensor by summing over $\theta$.
\end{definition}
\begin{theorem}[Exact delta-exterior completion]
    \label{thm:peps_labelled_open_delta}
    \lean{TNLean.PEPS.DependentBondNetwork.network_deltaCompletedTensor}
    \notready
    \uses{def:peps_labelled_open_coefficient}
    Fix a virtual label $v_0$. Replace every exterior site by a delta tensor
    fixing each incidence opposite $R$ to its prescribed boundary label,
    and each other exterior incidence to $v_0$. Retain the internal bond
    matrices and insert identity matrices on all other bonds. The closed
    contraction equals the open coefficient exactly, with no scalar factor.
\end{theorem}
\begin{proof}
    Split the site assignments into region and exterior assignments.
    The exterior deltas leave one assignment. Identity matrices on crossing
    bonds impose the boundary constraints, while those on exterior bonds
    contribute one. The remaining sum is the defining open coefficient.
\end{proof}
